% ! TeX root = ../../main.tex

% Observability for CORRECTNESS (not just performance) - ORCA as substrate Dynamic instrumentation
% (eBPF/DynInst) - claim as future work Generality within BSP: Charm++, MPI+X, ML 3D/4D parallelism
% all covered Out of scope: pure dataflow (Spark/Dask) - different causality model

\section{Related Work}\label{sec:orca:related}


\tabOrcaRelated{}

\parahead{ML Training Observability}
Recent ML training systems build application-specific pipelines spanning
straggler
localization~\cite{Wu2025:GREYHOUND,Yao2025:Holmes,Deng2025:Mycroft},
fault diagnosis~\cite{Deng2025:Minder,Dong2025:Aegis}, and correctness
enforcement~\cite{Jiang2025:TrainCheck}.
% % -- one paragraph version -- \parahead{ML Training Observability} Recent ML training systems
% build application-specific pipelines spanning straggler
% localization~\cite{Wu2025:GREYHOUND,Yao2025:Holmes,Deng2025:Mycroft}, fault
% diagnosis~\cite{Deng2025:Minder,Dong2025:Aegis}, and correctness
% enforcement~\cite{Jiang2025:TrainCheck}. ORCA generalizes the pattern.
% Greyhound~\cite{Wu2025:GREYHOUND} implements a complete feedback loop with NCCL event collection,
% microbenchmark validation, and online rebalancing. This is a fixed-function version of what ORCA
% expresses as composable operators. Holmes~\cite{Yao2025:Holmes} uses graph search over collective
% telemetry to find anomalies. ORCA's aligned collection makes stragglers surface through simple
% grouped aggregations, without requiring complex detection algorithms.
% Mycroft~\cite{Deng2025:Mycroft} buffers fine-grained NCCL telemetry locally and flushes to
% Kafka~\cite{Kreps2011:Kafka} on trigger, but notes the scalability limits of central logging.
% ORCA's tree-based overlay scales out naturally. TrainCheck~\cite{Jiang2025:TrainCheck} enforces
% training correctness via offline invariant generation and online checks. ORCA's timestep-aligned
% columnar analytics accelerate both stages. Minder~\cite{Deng2025:Minder} and
% Aegis~\cite{Dong2025:Aegis} operate at the cluster infrastructure level, using hardware counters
% and metrics to identify faulty nodes across workloads. This is a complementary scope to ORCA,
% which provides high-fidelity visibility into a single workload but does not yet support
% cross-workload diagnosis in multi-tenant environments.
ORCA generalizes the pattern.
Holmes~\cite{Yao2025:Holmes} uses graph search over collective telemetry to
find anomalies.
ORCA's aligned collection makes stragglers surface through
simple grouped aggregations, no complex detection needed.
Mycroft~\cite{Deng2025:Mycroft} buffers fine-grained NCCL~\cite{:NCCL} telemetry locally
and flushes to Kafka~\cite{Kreps2011:Kafka} on trigger, but notes the scalability limits of central
logging.
ORCA's tree-based overlay scales out naturally.
TrainCheck~\cite{Jiang2025:TrainCheck} enforces training correctness via
offline invariant generation and online checks.
ORCA's timestep-aligned
columnar analytics accelerate both stages.
Minder~\cite{Deng2025:Minder} and
Aegis~\cite{Dong2025:Aegis} operate at the cluster infrastructure level,
using hardware counters and metrics to identify faulty nodes across
workloads.
This targets a fundamentally different scope: identifying bad
hardware across many workloads treated as black boxes.
ORCA provides
high-fidelity visibility into a single workload but does not yet support
cross-workload diagnosis in multi-tenant environments.
\Cref{tab:orca:related-work} summarizes the diagnostic mechanisms implemented in these systems and
their ORCA equivalents.

Among these systems, Greyhound~\cite{Wu2025:GREYHOUND} is the closest to
ORCA in scope, implementing a complete feedback loop: iteration times are
tracked, NCCL events are collected for localization when outliers are
detected, microbenchmarks are dispatched to validate stragglers, and online
tuning rebalances work in response.
This is a fixed-function version of what ORCA expresses as composable operators.
ORCA's timestep dataframe abstraction preserves fine-grained trace events,
enabling disambiguation of stragglers without needing microbenchmarks for validation.
Greyhound's control semantics require pausing the application,
while ORCA implements an asynchronous control plane with atomicity and timestep-consistency
guarantees.
Additionally, ORCA's on-fabric broadcast tree enables faster message dissemination than a central
controller.
% arbitrary collectives to inject probes, while ORCA propagates updates asynchronously at step
% boundaries and uses on-fabric broadcast for faster dissemination than a central controller.

\parahead{Tool Infrastructure and Tracing}
MRNet~\cite{Roth2003:MRNet} introduced TBONs for scalable tool infrastructure, but its compiled C++
filter model has not evolved with modern analytics.
Fay~\cite{Erlin2011:Fay} and Snicket~\cite{Berg2021:Snicket} compile declarative queries into
distributed probes that filter and aggregate at source.
ORCA combines both: a TBON with Arrow interchange, declarative SQL, and BSP-aligned control.
Dynamic instrumentation via eBPF~\cite{:eBPF} and DynInst~\cite{Berna2011:DynInst} can extend reach
into runtime and kernel layers.
Sketch-based telemetry reduction~\cite{Sriva2024:Circa} aligns with ORCA's in-situ aggregation
model.

\parahead{Unified Analytics}
Arrow is emerging as common interchange for tabular HPC data ---TileDB~\cite{:TileDB} and
arrow-zarr~\cite{2026:arrow-zarr} bridge array stores and SQL engines~\cite{Lamb2024:DataFusion}.
SmartNICs are being explored as execution targets for Arrow-based data
pipelines~\cite{Ulmer2023:ExtendingComposableDatab,Liu2022:ProcessingParticleData}--- a declarative
engine like ORCA's could offload plan fragments to such devices transparently.
In-situ indexing systems like DeltaFS~\cite{Zheng2018:DeltaFS} and CARP~\cite{Jain2024:CARP} deploy
shufflers for embedded in-situ hash and range-partitioning.
ORCA's TBON could accommodate shuffling as an intra-tier operation.
As these ecosystems converge, common infrastructure can subsume analytics workflows ranging from
observability to checkpoint indexing to scientific queries.

% -- Details -- \parahead{ML Training Observability} Recent ML training systems address pathologies
% from straggler
% localization~\cite{Wu2025:GREYHOUND,Cui2025:XPUTimer,Yao2025:Holmes,Deng2025:Mycroft,Guan2025:PerfTracker,Chen2025:Reveal}
% to fault diagnosis~\cite{Deng2025:Minder,Dong2025:Aegis} to training
% correctness~\cite{Jiang2025:TrainCheck}, each building application-specific collection and
% analysis pipelines. % Megatron (NVIDIA), DeepSpeed (Microsoft) This is conceptual generalization.
% ORCA's interfaces will need to be extended to support 3D/4D parallelism, but the concepts apply.


% Greyhound implements a complete feedback loop: iteration times are tracked, NCCL events are
% collected for localization if outliers are detected, microbenchmarks are dispatched to validate
% stragglers, and online tuning interventions rebalance work in response. This is similar to a fixed
% function version of ORCA,

% 1. ORCA can materialize fine-grained statistics such as collective distribution histograms
% in-situ, more principled way to detect shifting straggler patterns in dynamic workloads like AMR
% codes. 2. In-situ aggregation and on-fabric transport allows trace events to be preserved for
% straggler attribution without necessarily requiring additional validation 3. Their control plane
% blocks at arbitrary collectives. ORCA's control plane asynchronous updates at step boundaries
% without needing to block. And on-fabric broadcast tree enables faster broadcast vs central
% controller.

% % Mycroft explicitly says that they don't currently scale to 10k+ GPUs Their eval is for 32 GPUs,
% % deployments at 128 GPUs or so Mycroft~\cite{Deng2025:Mycroft} uses fine-grained NCCL telemetry
% buffered locally and flushed to Kafka if heuristics trigger, but notes the scalability limits of
% central logging. ORCA's tree-based overlay enables scaling out.

% TrainCheck studies enforcing correctness in ML training via invariants. They use offline trace
% analysis for invariant generation, and online enforcement. ORCA emits timestep-aligned and
% partitioned Parquet, accelerating both offline analysis for invariant generation and in-situ
% enforcement.

% Holmes implements a graph search to analyze collective telemetry for anomalies. ORCA aligns
% collected dataframes, enabling simple GROUPBY's to identify stragglers.

% Reveal uses unsupervised learning on time-windowed metrics: ORCA can collect both metrics and
% trace events, but the latter is the source of truth.

% Minder and Aegis use cluster-level metrics and counters to identify faulty hardware. While these
% approaches are probalistic, they are more suited at cloud provider/infra level where workloads
% must be treated as black boxes. ORCA enables high-fidelity visibility into a single workload, but
% does not yet enable cross-workload visibility in multi-tenant environments.



% Arrays point: abandoned More broadly, scientific computing's array-native ecosystems
% (xarray~\cite{Hoyer2017:xarray}, Zarr~\cite{:Zarr}, HDF5~\cite{Folk2011:HDF5}) and relational
% analytics (Arrow, Parquet, SQL) are converging---TileDB~\cite{:TileDB} integrates array storage
% with database semantics, arrow-zarr~\cite{2026:arrow-zarr} bridges Zarr to DataFusion.



% \parahead{Tree-Based Overlay Networks} MRNet~\cite{Roth2003:MRNet} pioneered TBONs for scalable
% tool infrastructure, but its programming model—compiled C++ filter kernels over fixed packet
% formats—has not evolved with the data analytics ecosystem. ORCA revisits the TBON for modern
% infrastructure: receiver-driven RDMA transport, Arrow dataframes as interchange, declarative SQL,
% and timestep-consistent control.

% \parahead{Dynamic Instrumentation} ORCA's control plane generalizes naturally to dynamic
% instrumentation via eBPF~\cite{:eBPF} and DynInst~\cite{Berna2011:DynInst}, extending reach from
% application-level probes into runtime and kernel layers. eGPU~\cite{Yang2025:eGPU} extends dynamic
% instrumentation to GPUs. Deployment concerns around eBPF capability and privilege tradeoffs remain
% under active discussion ~\cite{Corbe2019:eBPFcap}. % in the Linux security community

% \parahead{Tracing and Observability} Fay~\cite{Erlin2011:Fay} and Snicket~\cite{Berg2021:Snicket}
% pioneered query-driven tracing, compiling declarative queries into distributed probes that filter
% and aggregate at the source rather than deferring to post-hoc analysis. ORCA adopts this paradigm
% for BSP workloads, adding BSP-aligned control, causal partitioning, and an Arrow-native on-fabric
% TBON. Sketch-based approaches to reducing telemetry volume via approximation are aligned with
% ORCA~\cite{Sriva2024:Circa}. eBPF-based tracing and observability mechanisms are deployed in
% production at scale~\cite{Benso2024:NetEdita}. otel-arrow~\cite{2026:Otelarrow} is a project to
% bridge OpenTelemetry and Arrow. \parahead{Trace Analytics} IOMax~\cite{Yildi2023:IOMax} converts
% I/O traces to Parquet and constructs cached aggregate views for interactive queries. ORCA writes
% Parquet natively, and our post-hoc evaluation (\S\ref{sec:orca:eval-trace}) demonstrates
% interactive queries over comparable data volumes using lazy Polars out of the
% box---dictionary-encoded Parquet, predicate pushdown, no bespoke optimization. For larger
% datasets, lakehouse formats~\cite{:Iceberg,Armbr2020:Deltalake} paired with distributed SQL
% engines~\cite{Zahar2010:Spark,Sethi2019:Prestoa,:Daft} scale out naturally.


% \parahead{ML Training} \cite{Yao2025:Holmes} (HOLMES: communication irregularity localization)
% \cite{Wu2025:GREYHOUND} (GREYHOUND: fail-slow detection) \cite{Deng2025:Minder} (Minder: faulty
% machine detection) \cite{Cui2025:XPUTimer} (XPUTimer: anomaly diagnostics) \cite{Dong2025:Aegis}
% (Aegis: failure diagnosis) \cite{Jiang2025:TrainCheck} (TrainCheck: silent erorrs, like numerical
% stability) \cite{Guan2025:PerfTracker}: Fine-grained profiline and differential observability to
% summarize Pytorch runtime behavior in prod \cite{Gangi2024:RDMA}: Tuning RDMA for jitter
% \cite{Chen2025:Reveal}: DLT anomalies using hardware counters + unsupervised learning
% \cite{Deng2025:Mycroft}: tracks collective-induced dependencies, fixed fault-localization
% pipeline, need multiple layers to debug "py-spy, Flight Recorder, Mycroft) - Mycroft also cites
% Dapper, X-Trace, Pivot Tracing, Hindsight for causal analyses \cite{Wang2025:Hawkeye}: PFC
% causality based diagnosis of RDMA performance anomalies ABCD

% \parahead{Towards Unified Analytics} ADIOS~\cite{Eisen2025:ADIOS} enables in-situ analytics via a
% coupled MPI application but lacks Arrow-native representation, declarative planning, or
% BSP-aligned control. DeltaFS~\cite{Zheng2018:DeltaFS} and CARP~\cite{Jain2024:CARP} deploy
% shufflers for embedded in-situ hash and range-partitioning. ORCA's TBON could accommodate
% shuffling as an intra-tier operation. More broadly, scientific computing's array-native ecosystems
% (xarray~\cite{Hoyer2017:xarray}, Zarr~\cite{:Zarr}, HDF5~\cite{Folk2011:HDF5}) and relational
% analytics (Arrow, Parquet, SQL) are converging---TileDB~\cite{:TileDB} integrates array storage
% with database semantics, arrow-zarr~\cite{2026:arrow-zarr} bridges Zarr to DataFusion.
% SmartNICs~\cite{Ulmer2023:ExtendingComposableDatab,Liu2022:ProcessingParticleData} offer
% Arrow-native acceleration for composable pipelines. ORCA's shared foundation positions it to
% benefit as these interfaces mature.

% \parahead{Streaming Analytics} ADIOS~\cite{Eisen2025:ADIOS} enables in-situ analytics via a
% coupled analytics plane running as a separate MPI application, whereas ORCA's analytics plane is
% fully asynchronous. ADIOS also lacks Arrow-native data representation, declarative query planning,
% operator pushdown, or BSP-aligned control semantics. DeltaFS~\cite{Zheng2018:DeltaFS} and
% CARP~\cite{Jain2024:CARP} deploy fixed-function pipelines for embedded indexing of checkpoint data
% using shuffle-based partitioning. ORCA's TBON topology and query planner could accommodate such
% workflows as intra-tier operations, enabling Dask-~\cite{Rockl2015:Dask} or
% Spark-style~\cite{Zahar2010:Spark} streaming analytics but in-situ.

% \parahead{Scientific Analytics} Scientific computing has historically centered on array-native
% ecosystems---libraries such as xarray~\cite{Hoyer2017:xarray} and Dask~\cite{Rockl2015:Dask} for
% labeled, chunked N-D computation and storage formats such as Zarr~\cite{:Zarr} and
% HDF5~\cite{Folk2011:HDF5}---while relational analytics evolved separately around SQL and columnar
% execution. Prior work on declarative interfaces for scientific data includes
% SciQL~\cite{Kerst2011:SciQL} for arrays, MeshSQL~\cite{Lee2004:MeshSQL} for region-centric mesh
% queries, and more recent algebras over unstructured
% meshes~\cite{Rezae2014:UnstructuredMeshAlgebra}. Practical convergence is now emerging:
% TileDB~\cite{:TileDB} integrates array storage with database-style semantics, while projects like
% arrow-zarr~\cite{2026:arrow-zarr} bridge Zarr stores to Arrow and DataFusion. ORCA's shared
% foundation positions it to benefit as these interfaces mature.

% \parahead{Accelerators} There is growing interest in leveraging accelerators such as SmartNICs for
% composable data pipelines, using Arrow as a shared
% representation~\cite{Ulmer2023:ExtendingComposableDatab,Liu2022:ProcessingParticleData}. ORCA's
% Arrow-native TBON can be extended to incorporate such devices as execution targets for plan
% fragments.


% Prototype implementations


% ORCA is highly extensible. % Cite sketches and all Right now we only do dynamic aggregation
% workflows over static telemetry. But we want to enable dynamic instrumentation via eBPF and
% DynInst to enable truly dynamic workflows.

% Long term, ideal observability remains the guiding principle. Having all telemetry be on tap, but
% cost is paid only when necessary. Currently, ORCA only gets to have the data after the trace
% collector such as Kineto has generated it. ORCA can drop/filter/aggregate it as close to the
% source, and it can mitigate the collection and transport overhead, but it can not undo the
% overhead Kineto has already incurred by recording the event. We think that control semantics allow
% deeper integration. Ideal observability is "collect only what the view necessitates, pay
% proportional cost. "

% Observability beyond performance and correctness. In situ visualization of meshes etc. is a thing.
% Can Science SQL enable it all via a unified platform?

% We think there is potential for modular interfaces for control and data. Private conversations
% suggest the challenge of managing IP constraints when cloud vendors are running customer workloads
% and performance problems occur. Standardized control and data interfaces can provide controlled
% visibility for services like ORCA to connect to and obtain more meaningful telemetry without
% sharing source code.

% Accelerators. Growing interest in using smart devices to manage telemetry volume. Also makes data
% pipelines complex. ORCA's Arrow + distributed query plannign allows these capabilities to be
% exposed in a portable way via the query planner, with users-side interfaces largely remaining
% unchanged.

% \section{Related Work}

% Analogy: high-frequency science sensors, like at SLAC/LHC, have a FPGA-based DAQ pipeline. It
% filters out 99\% of the data at source, while the rest is analyzed in post. ORCA attempts to be a
% similar software-defined thing.

% \textbf{In-situ}

% \cite{Zheng2018:DeltaFS} \cite{Godoy2020:ADIOS2} \cite{Larse2017:ALPINE}

% \textbf{Monitoring} \cite{Hintz2022:InfiniBandNetworkMonitoring}

% \textbf{Uncategorized} \cite{Wylie2025:15Years} \cite{Kerst2011:DataDeluge} \cite{Yokel2023:SOMA}
% -- "observability monitoring and in situ analytics in HPC" \cite{Eisen2025:ADIOS} -- streaming in
% HPC using ADIOS with SST "Sustainable Staging Transport" \cite{Knupf2012:Score-P} -- Score-P
% \cite{Schwa2021:LDMS} -- LDMS/Sandia? Idk \cite{Clasc:NSYS2PRVDetailedQuantitative} -- CLUSTER25,
% nsight to paraver, sqlite \cite{Bhate2019:Hatchet} -- callgraph index on pandas traces
% \cite{Bhate2024:Pipit} -- pandas based analytics on top of otf2, hpctoolkit etc
% \cite{Erlin2011:Fay} -- Theo recommended, distributed tracing, SOSP 11 \cite{Berg2021:Snicket} --
% query driven distributed tracing \cite{Devar2024:DFTracer} -- DFTracer, Hari et al
% \cite{Natar2010:TAUMRNet}

% \cite{Jiang2025:TrainCheck} \cite{Cui2025:XPUTimer} \cite{Deng2025:Minder} \cite{Dong2025:Aegis}
% \cite{Yao2025:Holmes} \cite{Wu2025:GREYHOUND} \cite{Xiong2024:SuperBench} \cite{Liang2025:Lumos}
% \cite{Yang2025:eGPU} -- extending eBPF to GPUs \cite{Jiang2024:MegaScale} -- scaling LLM training
% to 10k GPUS. has example of bootstrap problem (NCCL took 1047 seconds, they optimized it to 5s)

% \cite{Chen2025:Anomalies} -- "hardware-centric approach to

% \cite{Kerst2011:SciQL} -- SciQL \cite{Rezae2014:UnstructuredMeshAlgebra} -- SIGMOD
% \cite{Lee2004:MeshSQL} -- Mesh SQL \cite{:TileDB} -- TileDB

% Use Arrow: \cite{Liu2022:ProcessingParticleData} --- arrow for particle data and smartnics

% Do you really need to store all that telemetry? % X + MRNet ML Science SQL Blog post.
% \cite{Klein2024:YouReallyNeed}

% \cite{Kong2022:Collie} Benchmarks only validate performance for benchmark conditions. Real
% workloads generate unique conditions. They are very hard to recapture offline or in post.

% ADIOS \cite{Godoy2020:ADIOS2,Eisen2025:ADIOS} interesting architectural comparison wrt ORCA. Many
% higher-level differences: array-based formats vs Arrow-based dataframes. Analytics plane executes
% analytics baked into application, vs the whole TBON becomes a unified substrate where declarative
% analytics can be pushed down to point of origin. Also ORCA optimizes collection pipeline for TBON
% incast. More fundamentally, ORCA offers a completely asynchronous analytics plane, closer to a
% conventional distributed service. ADIOS analytics plane is another MPI application. ORCA
% aggregators are multi-threaded asynchronous loosely coupled.

% IOMax interesting. Manual ETL + optimizations + materialize entire dataframes in memory. ORCA
% shows that you can directly write partitioned Parquet, use dictionary-based columns, leverage
% engines with lazy execution to get great out of box performance. Recent versions of datafusion
% enable caching Parquet metadata to speed up subsequent queries without custom solutions.


% % ML Papers: 

% Aegis (Alibaba): no access to source code, instrument CCL % - Pull local stats if timeout detected
% (timeout ~ only fail-slow) % - If collectives started == n, it's a communication failure. If it's
% n - 1, it's a compute failure.

% Mycroft: fairly ORCA-shaped - Local logs/circular buffers of intra-collective events (fine-grained
% intra-NCCL flows) - Sampling based trigger: hoping it will catch stalls (only 10 ranks)
